\documentclass[aspectratio=169,11pt]{beamer}
\usetheme{TML}

\aula{1}{O que é Aprendizado de Máquina?}

\newcommand{\nbdir}{aulas/01-o-que-e-ml/notebooks}

\begin{document}

\maketitle

% =====================================================================
\begin{frame}{Três perguntas}
  \begin{enumerate}
    \item Quanto deve valer este apartamento de 70\,m² em Pinheiros?
    \item Este cliente vai pagar o empréstimo que está pedindo?
    \item Que ``tipos'' de famílias existem nos dados de orçamento doméstico?
  \end{enumerate}
  \bigskip
  \begin{columns}[T]
    \begin{column}{0.47\textwidth}
      \textbf{O que têm em comum?}\\
      Todas podem ser respondidas \alert{aprendendo padrões a partir de dados}
      de casos parecidos.
    \end{column}
    \begin{column}{0.47\textwidth}
      \textbf{Em que diferem?}\\
      No que queremos obter: um \alert{número}, uma \alert{categoria} ou
      \alert{grupos} que ninguém rotulou antes.
    \end{column}
  \end{columns}
  \bigskip
  \centering\small\textcolor{gray}{Ao fim da aula você saberá dar nome a cada uma e formulá-las como tarefas de ML.}
\end{frame}

\begin{frame}{Roteiro da aula}
  \setbeamertemplate{section in toc}[sections numbered]
  \small\tableofcontents[hideallsubsections]
\end{frame}

% =====================================================================
\section{O que significa ``aprender''?}

\begin{frame}{Uma definição operacional}
  \begin{definicao}
    Um programa \alert{aprende} com a experiência $E$, em relação a uma tarefa
    $T$ e a uma medida de desempenho $P$, se o seu desempenho em $T$, medido
    por $P$, melhora com $E$. \hfill \parencite{mitchell1997}
  \end{definicao}
  \medskip
  \begin{exemplo}[title={Exemplo: concessão de crédito}]
    \begin{tabular}{@{}ll@{}}
      $T$ (tarefa) & prever se um novo cliente ficará inadimplente \\
      $E$ (experiência) & histórico de clientes passados e se pagaram ou não \\
      $P$ (desempenho) & proporção de clientes \emph{novos} classificados corretamente
    \end{tabular}
  \end{exemplo}
  \medskip
  Repare: $P$ é medido em casos \alert{novos}. É isso que separa aprender de decorar.
\end{frame}

\begin{frame}{A mesma ideia em linguagem que você já conhece}
  Supomos que existe uma relação entre atributos $\vx$ e um alvo $y$:
  \[ y = f(\vx) + \varepsilon \]
  \begin{columns}[T]
    \begin{column}{0.5\textwidth}
      \begin{itemize}
        \item $f$ é \alert{desconhecida}
        \item Temos dados $\data = \{(\vx_i, y_i)\}_{i=1}^{n}$
        \item \alert{Aprender} = usar $\data$ para construir $\hf$
        \item \alert{Usar} = prever $\hat y = \hf(\vx_{\text{novo}})$
      \end{itemize}
    \end{column}
    \begin{column}{0.46\textwidth}
      \begin{conexao}
        \small Se $f(\vx) = \vx\T\vbeta$ e escolhemos $\hvbeta$ por mínimos
        quadrados, isso \emph{é} a regressão linear. O MQO já é um algoritmo
        de aprendizado!
      \end{conexao}
    \end{column}
  \end{columns}
\end{frame}

\begin{frame}{Dicionário econometria $\leftrightarrow$ aprendizado de máquina}
  \centering
  \begin{tabular}{@{}lll@{}}
    \toprule
    Econometria & Aprendizado de máquina & Símbolo \\
    \midrule
    variável dependente     & alvo, rótulo (\emph{target}, \emph{label})   & $y$ \\
    variáveis explicativas  & atributos (\emph{features})                    & $\vx$ \\
    observação              & exemplo, instância                             & $(\vx_i, y_i)$ \\
    amostra                 & conjunto de treino                             & $\dtreino$ \\
    estimar                 & treinar, ajustar (\emph{fit})                  & $\data \mapsto \hf$ \\
    parâmetros, coeficientes& pesos (\emph{weights})                         & $\vbeta$, $\vw$ \\
    intercepto              & viés (\emph{bias})                             & $\beta_0$, $b$ \\
    valor ajustado          & previsão (\emph{prediction})                   & $\hat y$ \\
    \bottomrule
  \end{tabular}
  \medskip

  \small\textcolor{gray}{Cuidado: ``viés'' em ML às vezes é só o intercepto, não o viés de um estimador.}
\end{frame}

% =====================================================================
\section{Mudando de domínio}

\begin{frame}{O percurso de sempre: começamos pelo DGP}
  \begin{columns}[T]
    \begin{column}{0.5\textwidth}
      \textbf{Pergunta:} quanto um ano a mais de estudo aumenta o salário?
      \medskip

      \textbf{Dados:} PNAD Contínua, uma linha por trabalhador.
      \medskip

      \textbf{Teoria:} o capital humano sugere o processo gerador dos dados
      (\emph{DGP}) \parencite{mincer1974}:
      \small
      \[ \ln w_i = \beta_0 + \beta_1\,\text{educ}_i + \beta_2\,\text{exper}_i
                   + \beta_3\,\text{exper}_i^2 + \varepsilon_i \]
    \end{column}
    \begin{column}{0.46\textwidth}
      \centering\small
      \begin{tabular}{@{}cccc@{}}
        \toprule
        educ & exper & exper$^2$ & $\ln w$ \\
        (anos) & (anos) & & \\
        \midrule
        11 & 20 &  400 & 7,60 \\
        16 &  8 &   64 & 8,16 \\
         9 & 30 &  900 & 7,31 \\
        18 & 15 &  225 & 8,99 \\
        \bottomrule
      \end{tabular}
      \par\smallskip
      {\footnotesize\textcolor{gray}{valores ilustrativos}}
      \medskip
      \begin{intuicao}
        \small Cada coluna de $\mX$ tem \alert{nome, unidade e significado}.
      \end{intuicao}
    \end{column}
  \end{columns}
\end{frame}

\begin{frame}{Cada termo conta uma história}
  \[ \ln w_i = \beta_0 + \beta_1\,\text{educ}_i + \beta_2\,\text{exper}_i
               + \beta_3\,\text{exper}_i^2 + \varepsilon_i \]
  \begin{itemize}
    \item<2-> $\beta_1$: um ano a mais de estudo muda o salário em cerca de
          $100\,\beta_1$\%: o \alert{retorno da educação}
    \item<3-> $\beta_2 > 0$ e $\beta_3 < 0$: o salário cresce com a
          experiência, mas cada vez menos (\alert{concavidade})
    \item<4-> $\varepsilon_i$: o que não observamos (habilidade, rede de
          contatos\ldots). Se a habilidade se correlaciona com educ, $\hat\beta_1$
          tem \alert{viés de variável omitida}
  \end{itemize}
  \begin{onlyenv}<5>
  \begin{conexao}[title={O percurso que vocês conhecem}]
    \small teoria $\to$ DGP $\to$ estimar $\hvbeta$ $\to$ testar hipóteses
    $\to$ \alert{interpretar} cada coeficiente
  \end{conexao}
  \end{onlyenv}
\end{frame}

\begin{frame}{Por que esse percurso funciona?}
  Ele se apoia em três condições que quase nunca questionamos:
  \medskip
  \begin{enumerate}
    \item \textbf{Poucas colunas}: $p = 4$ variáveis, e o número de
          observações $n$ é muito maior que $p$.
    \item \textbf{Cada coluna significa algo sozinha}: ``educ'' é educação
          para todo mundo, e a ordem das colunas não importa.
    \item \textbf{A teoria sugere a forma de $f$}: linear em educ, quadrática
          em exper.
  \end{enumerate}
  \medskip
  \begin{atencao}
    \small Guarde essas três condições. Nos próximos slides vamos quebrá-las,
    uma a uma.
  \end{atencao}
\end{frame}

\begin{frame}{E se mudarmos de domínio?}
  \begin{columns}[c]
    \begin{column}{0.5\textwidth}
      \includegraphics[width=\linewidth]{a01_gato_cachorro}
    \end{column}
    \begin{column}{0.46\textwidth}
      \Large
      $\vx$ = os \alert{pixels} de uma imagem\\[6pt]
      $y \in \{\text{gato}, \text{cachorro}\}$
      \normalsize
      \bigskip

      \begin{uncoverenv}<2>
      \textbf{Pergunta para a turma:} escrevam o DGP.
      \[ y_i = \;?\]
      \end{uncoverenv}
    \end{column}
  \end{columns}
\end{frame}

\begin{frame}{O que o computador vê}
  \centering
  \includegraphics[width=0.54\linewidth]{a01_pixels}
  \par\smallskip\raggedright\small
  Cada pixel é um número de 0 (preto) a 255 (branco). A imagem $24\times24$
  vira \alert{uma linha} de $\mX$ com 576 colunas. Uma foto de
  $224\times224$ em cores (3 canais) tem $224 \times 224 \times 3 = 150\,528$.
  \begin{atencao}
    \small \textbf{Condição 1 quebrada.} De 4 colunas passamos a centenas de
    milhares, e $p$ pode ser maior que $n$.
  \end{atencao}
\end{frame}

\begin{frame}{Vizinhança: o significado está no arranjo}
  \centering
  \includegraphics[width=0.6\linewidth]{a01_deslocamento}
  \par\smallskip\raggedright\small
  O pixel $(12, 7)$, coluna 295 de $\mX$ (contando do zero), é \alert{olho}
  numa imagem e \alert{pelo} na outra. Mudaram 128 colunas e $y$ continua
  ``gato'': o que importa é o padrão que cada pixel forma com os
  \alert{vizinhos}.
  \begin{atencao}
    \small \textbf{Condição 2 quebrada.} Na PNAD, a coluna educ é sempre
    educação. Aqui, nenhuma coluna tem significado próprio.
  \end{atencao}
\end{frame}

\begin{frame}{Vamos tentar com o que sabemos}
  \begin{enumerate}
    \item<1-> \textbf{Logit nos pixels:}
          $P(\text{gato}\mid\vx) = \Lambda\big(\beta_0 + \sum_{j=1}^{150\,528}\beta_j x_j\big)$\\
          \small O que é $\beta_{295}$? ``O efeito do brilho do pixel 295 sobre a
          chance de ser gato''. E se o gato estiver 2 pixels para o lado?
          \normalsize
    \item<2-> \textbf{Deixar mais flexível com interações:} só os pares de pixels
          já dão $150\,528 \times 150\,527 / 2 \approx 11{,}3$ \alert{bilhões} de
          coeficientes.
    \item<3-> \textbf{Criar atributos à mão:} ``tem orelha pontuda?'',
          ``tem focinho?''\\
          \small Como escrever ``orelha'' em termos de pixels, em qualquer
          posição, tamanho, iluminação, pose e raça?\normalsize
  \end{enumerate}
  \begin{onlyenv}<4>
  \begin{atencao}
    \small \textbf{Condição 3 quebrada.} Não há teoria que nos dê a forma de
    $f$. Se bateu um certo desespero, é isso mesmo: o percurso do DGP
    \alert{não tem por onde começar}.
  \end{atencao}
  \end{onlyenv}
\end{frame}

\begin{frame}{Dois mundos de $\mX$}
  \centering\small
  \begin{tabular}{@{}>{\raggedright\arraybackslash}p{0.27\textwidth}
                    >{\raggedright\arraybackslash}p{0.3\textwidth}
                    >{\raggedright\arraybackslash}p{0.33\textwidth}@{}}
    \toprule
    & Salário (PNAD) & Imagem (gato ou cachorro) \\
    \midrule
    nº de colunas $p$              & 4                       & de centenas a centenas de milhares \\
    significado de uma coluna      & próprio (anos de estudo) & nenhum, isoladamente \\
    ordem das colunas              & irrelevante             & essencial (vizinhança) \\
    mesmo $y$, $\vx$ muito diferente & raro                  & sempre: posição, luz, escala, pose \\
    colunas vizinhas               & algumas correlacionadas & quase iguais (colinearidade extrema) \\
    de onde vem $f$                & da teoria               & não há teoria \\
    pergunta                       & $\hat\beta$: quanto vale estudar? & $\hat y$: é um gato? \\
    \bottomrule
  \end{tabular}
\end{frame}

\begin{frame}{É para isso que usamos aprendizado de máquina}
  \begin{columns}[T]
    \begin{column}{0.5\textwidth}
      \textbf{Invertemos o percurso:}
      \begin{itemize}\small
        \item não escrevemos $f$; damos ao algoritmo \alert{muitos exemplos}
              $(\vx_i, y_i)$ rotulados
        \item ele escolhe $\hf$ dentro de uma família \alert{muito flexível}
              de funções
        \item julgamos $\hf$ pelos acertos em \alert{imagens novas}, como na
              definição de Mitchell
      \end{itemize}
      \medskip
      \small Perdemos coeficientes interpretáveis; ganhamos previsões onde
      não há teoria para $f$.
    \end{column}
    \begin{column}{0.46\textwidth}
      \begin{exemplo}[title={E economistas com isso?}]
        \small Pixels de imagens de satélite servem para prever a pobreza de
        regiões com poucos dados de pesquisas domiciliares \parencite{jean2016}. Texto e
        imagens como dado: Aula~6.
      \end{exemplo}
      \begin{atencao}
        \small O ML não aposenta a econometria. Se a pergunta é o retorno da
        educação, o percurso do DGP continua sendo o caminho.
      \end{atencao}
    \end{column}
  \end{columns}
\end{frame}

% =====================================================================
\section{Estatística, econometria e ML}

\begin{frame}{Duas culturas}
  \begin{columns}[c]
    \begin{column}{0.6\textwidth}
      \includegraphics[width=\linewidth]{a01_duas_culturas}
    \end{column}
    \begin{column}{0.37\textwidth}
      \textcite{breiman2001cultures}:
      \begin{itemize}\small
        \item a estatística tradicional parte de um \alert{modelo para os dados}
              e se interessa pelos parâmetros (o salário);
        \item o ML parte de um \alert{algoritmo que prevê bem} e o julga pelo
              desempenho fora da amostra (o gato).
      \end{itemize}
      \small A matemática é a mesma; muda o \textbf{produto final}.
    \end{column}
  \end{columns}
\end{frame}

\begin{frame}{Problemas de $\hat\beta$ e problemas de $\hat y$}
  \begin{columns}[T]
    \begin{column}{0.48\textwidth}
      \textbf{\textcolor{tmlblue}{Problemas de $\hat\beta$}}\\[2pt]
      Queremos entender um \alert{mecanismo}.
      \begin{itemize}\small
        \item Qual o efeito de um ano a mais de estudo sobre o salário?
        \item O Bolsa Família afeta a frequência escolar?
      \end{itemize}
      \small Exige identificação: o que aconteceria \emph{se} mudássemos $x$?
    \end{column}
    \begin{column}{0.48\textwidth}
      \textbf{\textcolor{tmlaccent}{Problemas de $\hat y$}}\\[2pt]
      Queremos uma boa \alert{previsão}.
      \begin{itemize}\small
        \item Qual a probabilidade de este cliente ficar inadimplente?
        \item Qual será o IPCA do próximo mês?
      \end{itemize}
      \small Exige generalização: acertar em casos que ainda não vimos.
    \end{column}
  \end{columns}
  \bigskip
  \centering
  O aprendizado de máquina é feito para problemas de $\hat y$
  \parencite{mullainathan2017}.
\end{frame}

\begin{frame}{Muitos $\hat\beta$ diferentes, quase o mesmo $\hat y$}
  \centering
  \includegraphics[width=0.8\linewidth]{a01_prever_inferir}
  \par\raggedright\small
  1000 amostras de $y = x_1 + x_2 + \varepsilon$, com $\operatorname{corr}(x_1,x_2)=0{,}98$.
  O coeficiente de $x_1$ \alert{oscila muito}; a previsão em $x_0=(1,1)$ \alert{quase não muda}.
  \begin{atencao}
    \small Um modelo que prevê bem \textbf{não} garante coeficientes
    interpretáveis. Não leia os ``pesos'' de um modelo preditivo como efeitos causais.
  \end{atencao}
\end{frame}

\begin{frame}{Guarda-chuva ou dança da chuva?}
  \begin{columns}[T]
    \begin{column}{0.48\textwidth}
      \textbf{Levar o guarda-chuva?}\\
      Basta \alert{prever} se vai chover. Não importa \emph{por que} chove.
      \medskip

      \small Política: quais pacientes idosos provavelmente morrerão no ano
      seguinte, e para quem uma prótese de quadril traria pouco benefício?
      \emph{Previsão} orienta a decisão.
    \end{column}
    \begin{column}{0.48\textwidth}
      \textbf{Fazer a dança da chuva?}\\
      Precisamos saber se a dança \alert{causa} chuva. Prever não basta.
      \medskip

      \small Política: um programa de reforço escolar \emph{reduz} a evasão?
      Exige \emph{inferência causal}.
    \end{column}
  \end{columns}
  \bigskip
  \begin{conexao}
    \small Muitas decisões públicas são \emph{problemas de previsão}
    \parencite{kleinberg2015}. E mesmo estudos causais usam ML em etapas
    preditivas \parencite{chernozhukov2018,athey2019}, assunto da Aula 6.
  \end{conexao}
\end{frame}

% =====================================================================
\section{Tipos de aprendizado}

\begin{frame}{A pergunta-chave: temos o alvo $y$?}
  \centering
  \includegraphics[height=0.78\textheight]{a01_taxonomia}
\end{frame}

\begin{frame}{As três tarefas em uma imagem}
  \centering
  \includegraphics[width=\linewidth]{a01_tarefas}
  \par\raggedright\small
  Regressão e classificação \alert{aprendem com respostas conhecidas}.
  O agrupamento \alert{propõe uma estrutura}, e cabe a nós interpretá-la
  (e desconfiar dela).
\end{frame}

\begin{frame}{Formular é escolher}
  A mesma base de dados pode gerar tarefas diferentes:
  \medskip
  \begin{center}
  \begin{tabular}{@{}llp{0.38\textwidth}@{}}
    \toprule
    Alvo escolhido & Tarefa & O que se ganha ou perde \\
    \midrule
    renda em R\$ & regressão & usa toda a informação \\
    renda $>$ 2 salários mínimos? & classificação & mais simples de comunicar; descarta informação \\
    (nenhum alvo) & agrupamento & descobre perfis, sem garantia de que façam sentido \\
    \bottomrule
  \end{tabular}
  \end{center}
  \begin{intuicao}
    \small A formulação vem \textbf{antes} do algoritmo e depende da \alert{decisão}
    que a previsão vai apoiar. Errar aqui não se corrige depois com um modelo melhor.
  \end{intuicao}
\end{frame}

% =====================================================================
\section{Formulando problemas econômicos}

\begin{frame}{Ficha de formulação}
  \begin{columns}[T]
    \begin{column}{0.44\textwidth}
      \begin{enumerate}\small
        \item \textbf{Decisão}: quem usa a previsão e para quê?
        \item \textbf{Unidade}: o que é uma linha?
        \item \textbf{Alvo $y$}: o que prever?
        \item \textbf{Atributos $\vx$}: o que se sabe \alert{no momento da decisão}?
        \item \textbf{Tarefa}: regressão, classificação, agrupamento?
        \item \textbf{Métrica}: como saber se ficou bom?
      \end{enumerate}
    \end{column}
    \begin{column}{0.52\textwidth}
      \begin{exemplo}[title={Exemplo: crédito pessoal}]
        \small
        \textbf{Decisão:} aprovar ou não o pedido\\
        \textbf{Unidade:} um pedido de empréstimo\\
        \textbf{Alvo:} atraso $>$ 90 dias em 12 meses (sim/não)\\
        \textbf{Atributos:} renda, idade, dívida/renda, histórico\\
        \textbf{Tarefa:} classificação\\
        \textbf{Métrica:} erros custam diferente: negar um bom pagador
        $\neq$ aprovar um mau pagador (Aula 3)
      \end{exemplo}
    \end{column}
  \end{columns}
\end{frame}

\begin{frame}{A armadilha do item 4}
  \begin{atencao}
    Usar um atributo que \textbf{só existe depois} do que queremos prever.
  \end{atencao}
  \medskip
  \begin{itemize}
    \item Prever inadimplência usando ``número de cobranças enviadas''
    \item Prever o PIB do trimestre usando a \emph{revisão final} do próprio PIB
    \item Prever evasão escolar usando ``data de cancelamento da matrícula''
  \end{itemize}
  \medskip
  O modelo parece excelente nos dados históricos e é \alert{inútil na prática}.
  Esse é o \emph{vazamento de dados}, tema central da Aula 3.
\end{frame}

% =====================================================================
\section{Fluxo de trabalho em Python}

\begin{frame}{O fluxo de um projeto}
  \centering
  \includegraphics[height=0.52\textheight]{a01_fluxo}
  \par\medskip\raggedright
  \begin{itemize}\small
    \item \textcolor{tmlaccent}{\textbf{Laranja}}: separar e avaliar sem se enganar (Aula 3)
    \item \textcolor{tmlblue}{\textbf{Azul}}: como os modelos são treinados por dentro (Aulas 4 e 5)
    \item O ciclo \emph{avaliar $\to$ ajustar} se repete, mas o conjunto de teste
          só é usado \alert{uma vez}, no final
  \end{itemize}
\end{frame}

\begin{frame}[fragile]{scikit-learn: uma interface para todos os modelos}
  \begin{codigo}
    from sklearn.model_selection import train_test_split
    from sklearn.linear_model import LinearRegression
    from sklearn.metrics import mean_squared_error

    X_tr, X_te, y_tr, y_te = train_test_split(X, y, test_size=0.2, random_state=0)

    modelo = LinearRegression()      # 1. escolher o modelo
    modelo.fit(X_tr, y_tr)           # 2. aprender com o treino
    y_hat = modelo.predict(X_te)     # 3. prever casos não vistos
    mean_squared_error(y_te, y_hat) ** 0.5   # 4. avaliar fora da amostra
  \end{codigo}
  \begin{intuicao}
    \small Todo modelo do \pkg{scikit-learn} \parencite{pedregosa2011} tem
    \py{fit} e \py{predict}. Trocar a regressão linear por uma floresta aleatória
    é trocar \textbf{uma linha}.
  \end{intuicao}
\end{frame}

\begin{frame}[fragile]{Sempre compare com uma referência ingênua}
  \begin{columns}[T]
    \begin{column}{0.5\textwidth}
      \begin{codigo}
        from sklearn.dummy import DummyRegressor

        ingenuo = DummyRegressor(strategy="mean")
        ingenuo.fit(X_tr, y_tr)
        # prevê sempre a média do treino
      \end{codigo}
    \end{column}
    \begin{column}{0.46\textwidth}
      \begin{itemize}\small
        \item Regressão: prever sempre a \alert{média}
        \item Classificação: prever sempre a \alert{classe mais comum}
        \item Séries temporais: prever o \alert{último valor}
      \end{itemize}
    \end{column}
  \end{columns}
  \medskip
  \begin{atencao}
    \small Se 95\% dos clientes pagam, um ``modelo'' que diz \emph{todos pagam}
    acerta 95\%. Um número de desempenho só tem significado \textbf{comparado} a
    uma referência.
  \end{atencao}
\end{frame}

% =====================================================================
\section{Mãos à obra}

\demo[\nbdir/01_demo.ipynb]{Do problema ao primeiro modelo:\\ preços de imóveis com scikit-learn}

\lab[\nbdir/01_lab.ipynb]{Crédito e perfis de consumo:\\ classificação e agrupamento}

\begin{frame}{Para fixar}
  \begin{itemize}\small
    \item Aprender = melhorar o desempenho em \alert{casos novos} com a experiência
    \item Sem colunas com significado próprio e sem teoria para $f$ (imagens,
          texto), o percurso do DGP trava: é aí que o ML entra
    \item $\hat y$ e $\hat\beta$ respondem a perguntas diferentes; bons $\hat y$ não
          garantem $\hat\beta$ interpretáveis
    \item Supervisionado (regressão, classificação) $\times$ não supervisionado
          (agrupamento, redução de dimensão)
    \item Formular vem antes de modelar: decisão, unidade, alvo, atributos
          \alert{disponíveis}, tarefa, métrica
    \item \py{fit} / \py{predict} e uma referência ingênua, sempre
  \end{itemize}
  \medskip
  \small\textbf{Para aprofundar:} \textcite{mullainathan2017}; \textcite[cap.~2]{james2021};
  \textcite{breiman2001cultures}.
  \smallskip

  \textbf{Próxima aula:} Python, NumPy e JAX, as ferramentas para abrir os modelos por dentro.
\end{frame}

\begin{frame}{Referências}
  \printbibliography[heading=none]
\end{frame}

\end{document}
